% !TEX root = ../main.tex
\chapter*{Requests about these notes}\label{ch:requests}
\addcontentsline{toc}{chapter}{Requests about these notes}
\markboth{Requests about these notes}{}
\begin{paracol}{2}

Your margin notes are of two kinds: notes on the \emph{content} (a doubt, an idea), which stay next to the
text they refer to, and requests about the \emph{notes themselves} (layout, boxes, conventions), which are
collected here with what was done about them.

\paragraph{How the margin notes are marked.}
\begin{itemize}
  \item \verb|\note{...}|: a \emph{new} note, not yet read. ``Read my notes'' means these, and only these.
  \item A note that has been answered becomes one box with the question on top and a dark yellow
        \textbf{ANSWER} bar below it (\verb|\begin{answer}{question}[topic]| \dots\ \verb|\end{answer}|).
  \item \verb|\note*{...}|: a request that has been carried out (green DONE tag).
  \item \verb|\note+{...}|: a comment that has been read and needs no answer (blue SEEN tag).
  \item \verb|\feedback{...}|: a comment on one of your own solutions (indigo).
\end{itemize}

\section*{The list}
No requests yet.

\flushside
\end{paracol}
